\documentclass[10pt,a4paper]{amsart}
\usepackage[margin=19mm]{geometry}
\usepackage{amsmath,amssymb,mathtools}
\usepackage[scr=rsfs,cal=euler]{mathalfa}
\usepackage{xfrac}
\usepackage[unicode,hidelinks,bookmarks=false]{hyperref}
\newcommand{\ie}{i.e.\ }
\newcommand{\cf}{cf.\ }
\newcommand{\resp}{respectively\ }
\newcommand{\Subsection}{\subsection}
\providecommand{\nobreakdash}{\nobreak\hbox{-}}
\setlength{\emergencystretch}{2em}
\multlinegap0pt
\usepackage{bm}
\usepackage{bbm}
\usepackage{dsfont}
\usepackage{xspace}
\usepackage{cancel}
\usepackage{mathtools}
\usepackage{amsthm}
\makeatletter
\@ifundefined{theorem}{\newtheorem{theorem}{Theorem}}{}
\@ifundefined{lemma}{\newtheorem{lemma}[theorem]{Lemma}}{}
\makeatother
\let\g=\gamma \let\h=\eta    \let\k=\kappa  \let\l=\lambda
\newcommand{\E}{\mathds{E}}
\newcommand{\car}{\mathds{1}}
\newcommand{\scalar}[2]{\left\langle #1, #2 \right\rangle}
\title[A universal cutoff phenomenon for mean-field exchange models]{A universal cutoff phenomenon for mean-field exchange models}
\author{Pietro Caputo, Matteo Quattropani, Federico Sau}
\date{Source: arXiv:2506.12816v1 (CC BY 4.0)}
\begin{document}
\maketitle

\noindent\emph{Source and licence.} A universal cutoff phenomenon for mean-field exchange models. Version arXiv:2506.12816v1 (CC BY 4.0), distributed under the Creative Commons Attribution 4.0 International licence, as recorded in the arXiv licence metadata for this exact version and shown on the abstract page https://arxiv.org/abs/2506.12816v1. Full rights notice: licenses/arxiv-2506.12816-CC-BY-4.0.txt.\par\smallskip

\noindent\textit{Editorial scope.} Counted unit: the Lemma whose source label is \texttt{lemma:eta*B} in the article (source0.tex lines 789--794, source label \texttt{lemma:eta*B}), delivered together with the article's own complete original proof of it (source0.tex lines 832--901, one \texttt{proof} environment of 70 lines). No mathematics is written, repaired or re-derived: statement and proof are the article's own text line for line. The proof is detached from the statement in the source: the article states the result at lines 789--794 and prints its proof at lines 832--901, and the proof environment's own header names the statement, so the association is the article's own. Required context retained uncounted: eq:def-eta+g (lines 406--408); eq:dual (lines 489--491); eq:def-eta-star (lines 729--731). Every definition, display and label of those lines is retained unchanged. Omitted from this package and not redistributed: 1--405, 409--488, 492--728, 732--788, 795--831, 902--1116. Nothing inside a retained excerpt is omitted or altered: every retained body is delivered exactly as the article's lines are written, so no omission receipt of any kind accompanies this package. Citations: the retained excerpts carry no citation command at all, so no bibliography entry of the article is transcribed and no reference is redistributed. Reference adaptation: none was needed and none was made; every reference inside the delivered unit resolves inside it, so no unresolved reference or citation is produced. Printed slips kept literal: no printed word, symbol, hypothesis or number of the counted statement or of its proof was altered. Layout and class: the article's own class is \texttt{amsart} with packages including babel, comment, amsmath, dsfont, mathtools, amssymb, color, hyperref, appendix; the wrapper is the lane's pinned \texttt{amsart} editorial wrapper at 10pt on A4, which restates the theorem-like environments the retained text needs and adds the article's own macro definitions that the retained text needs (\texttt{\textbackslash E}, \texttt{\textbackslash car}, \texttt{\textbackslash g}, \texttt{\textbackslash scalar}), with extra wrapper packages (bm, bbm, dsfont, xspace, cancel, mathtools). The delivered numbering is the unit's own. Assets: no retained span contains a figure environment, a graphics-inclusion command, a PDF-page inclusion, a table or a TikZ picture; no image file is copied into this package, the package declares no asset dependency and no image file is redistributed with it. Licence: version arXiv:2506.12816v1 of this article is distributed under the Creative Commons Attribution 4.0 International licence, as recorded in the arXiv licence metadata for this exact version and shown on the abstract page https://arxiv.org/abs/2506.12816v1. A full rights notice is delivered at licenses/arxiv-2506.12816-CC-BY-4.0.txt. Characters: every retained excerpt is pure ASCII, so the delivered engine, XeLaTeX with its default Latin Modern text font, reports no missing character.
\begin{equation}\label{eq:def-eta+g}
\eta_t^{+,\g}(x)=\sum_{\zeta\in A_t(x)} |\zeta|\car_{n|\zeta|\ge e^{\psi}}\,.
\end{equation}

\begin{equation}\label{eq:dual}
	\scalar{\eta_0}{\xi_t}\overset{d}{=}\scalar{\eta_t}{\xi_0} ,
\end{equation}

\begin{equation}\label{eq:def-eta-star}
		\eta_{t}^*(x)=\eta_{t}^-(x)\car_{n\eta_{t}^-(x)< e^{3\psi}}\,,\qquad \widehat \eta_t = \eta_{t}-\eta_{t}^{*}\,.
	\end{equation}

	\begin{lemma}\label{lemma:eta*B}
	All three models satisfy
		\begin{equation}\label{eq:topro}
			\lim_{n\to\infty}\big(\E[\|\widehat\eta_{t}\|_1]-\E[\|\eta^+_{t}\|_1]\big)=0\,.
		\end{equation}
	\end{lemma}

\begin{proof}[Proof of Lemma \ref{lemma:eta*B}]	Recall the definitions \eqref{eq:def-eta+g} and \eqref{eq:def-eta-star}, and define
	\begin{equation}
		B_t=\{x:\, n\eta_{t}^-(x)\ge e^{3\psi} \}\,.
	\end{equation}
	In this way \eqref{eq:topro} reduces to 
		\begin{equation}\label{eq:topros0}
			\lim_{n\to\infty}\E\left[\eta_{t}^-(B_t)\right]=0\,,
		\end{equation}
		where we use the notation 
		\begin{equation}\label{eq:zeta-B}
				\eta_t^-(B_t)=\sum_{x\in B_t}\eta^-_t(x) = \sum_{x\in B_t}\sum_{\zeta\in A_t(x)}|\zeta|\car_{n|\zeta|<e^\psi}
				\,.
				\end{equation}
We will actually prove that \eqref{eq:topros0} holds uniformly in $t\ge 0$, i.e.\
\begin{equation}\label{eq:topros}
			\lim_{n\to\infty}\sup_{t\ge 0}\E\left[\eta_{t}^-(B_t)\right]=0\,,
		\end{equation}
Notice that
		\begin{align}
				\eta_t^-(B_t)&=\sum_{x}\car_{n\eta_t^-(x)\ge e^{3\psi}}\sum_{\zeta\in A_t(x)}|\zeta|\car_{n|\zeta|< e^{\psi}}\\
				&\le ne^{-3\psi}\sum_{x}\eta_t^-(x)\sum_{\zeta\in A_t(x)}|\zeta|\car_{n|\zeta|< e^{\psi}}
				= ne^{-3\psi}\sum_{x}\eta_t^-(x)^2\,.
		\label{eq:zeta-B2}\end{align}

		We now estimate  the expectation of the sum in the right hand side of \eqref{eq:zeta-B2}. One has 
		\begin{align}\label{eq:Ed-End}
	\eta_t^-(x)^2=\sum_{\zeta,\zeta'\in A_t(x)}|\zeta|\,|\zeta'|\car_{n|\zeta|<{e^\psi}}\car_{n|\zeta'|<{e^\psi}}\le E_t^{\rm d}(x)+E_t^{\rm nd}(x)\,,
			\end{align}
	where we define 
		\begin{align}
	E_t^{\rm d}(x)=\sum_{\zeta\in A_t(x)}|\zeta|^2\car_{n|\zeta|<{e^\psi}}
	\,,\qquad 
	E_t^{\rm nd}(x) = \sum_{\zeta,\zeta'\in A_t(x)}|\zeta|\,|\zeta'|\car_{\zeta\neq \zeta'}\,,
			\end{align}
		The expectation of the sum of the diagonal terms $E_t^{\rm d}(x)$ can be bounded as
		\begin{equation}\label{eq:zeta-B3}
			\sum_{x}\E[E_t^{\rm d}(x)]\le \frac{e^\psi}{n}\sum_{x}\E\left[ \sum_{\zeta\in A_t(x)}|\zeta|\car_{|\zeta|<\frac{e^\psi}{n}}\right]=\frac{e^\psi}{n}\E[\|\eta_t^-\|_1]\le \frac{e^\psi}{n}\,.
		\end{equation}
			For the non-diagonal terms $E_t^{\rm nd}(x)$, we argue as follows.  Let $(u,v)$ denote the random pair of particles to be updated at time $t$ and let $X$ denote the associated redistribution variable. If $x\notin (u,v)$ then nothing happens at $x$ and therefore $A_t(x) = A_{t-1}(x)$ in this case. It follows that
			 \begin{equation}\label{eq:xnot}
			\E\left[\car_{x\notin (u,v)}E_t^{\rm nd}(x)\right] = \left(1-\frac2n\right)\E\left[E^{\rm nd}_{t-1}(x)\right]\,.
					\end{equation}
					If $(u,v) = (x,y)$, then any pile $\zeta\in A_t(x)$ has size $X|\zeta_0|$ where $\zeta_0$ is a uniquely determined pile, the parent, that was either in $A_{t-1}(x)$ or $A_{t-1}(y)$. This is true for the SRM. 	We will stick with the case of the SRM for the sake of simplicity, and will comment on the minor modifications needed in the other two models later on. We note that in any case $\zeta\neq\zeta'$ implies $\zeta_0\neq\zeta_0'$. Therefore, for a fixed pair $(x,y)$,
			 \begin{align}\label{eq:xin}
			\E\left[\car_{(x,y)= (u,v)}E_t^{\rm nd}(x)\right]& =\frac{\E[X^2]}{n(n-1)}\E\left[\sum_{\zeta,\zeta'\in A_{t-1}(x)\cup A_{t-1}(y)}|\zeta|\,|\zeta'|\car_{\zeta\neq \zeta'}
			\right]\,. 
					\end{align}
					The last expectation above can be rewritten as 
					 \begin{align}\label{eq:xin2}
					 \E\left[\sum_{\zeta,\zeta'\in A_{t-1}(x)\cup A_{t-1}(y)}|\zeta|\,|\zeta'|\car_{\zeta\neq \zeta'}
			\right]= \E[E_{t-1}^{\rm nd}(x)] + \E[E_{t-1}^{\rm nd}(y)]+2\E[\eta_{t-1}(x)\eta_{t-1}(y)]\,.
				\end{align}	
			The case $(u,v)=(y,x)$ is analogous, with $X$ replaced by $1-X$. Since $\E[X^2]=\E[(1-X)^2]$, the expressions \eqref{eq:xnot} and \eqref{eq:xin} imply
		\begin{align}
			\sum_{x}\E[E_t^{\rm nd}(x)]&\le  \left(1-\frac2n\right)\sum_{x}\E\left[E^{\rm nd}_{t-1}(x)\right] +\\& \quad + \frac{2\E[X^2]}{n(n-1)}{\sum_{x,y\,:\,x\neq y}}\left(\E[E_{t-1}^{\rm nd}(x)] + \E[E_{t-1}^{\rm nd}(y)]+2\E[\eta_{t-1}(x)\eta_{t-1}(y)]\right)\\ & 
			\leq \left(1-\frac2n + \frac{4\E[X^2]}n\right)\sum_{x}\E\left[E^{\rm nd}_{t-1}(x)\right] + \frac{4\E[X^2]}{n(n-1)}\,\E\left[\|\eta_{t-1}\|_1^2\right]\,.
					\end{align}
		The above estimate has been derived for SRM. Let us remark that the same expression holds for SEM and GAM with the only modification that 	the last term in the right hand side has $\E[X^2]$ replaced by $\E[X(1-X)]$. This follows from symmetry of $X$ and the fact that 
		in GAM and SEM one has $\zeta=X|\zeta_0|$ if $\zeta_0\in A_{t-1}(x)$ and $\zeta=(1-X)|\zeta_0|$ if $\zeta_0\in A_{t-1}(y)$.
Moreover, for SRM and GAM		one has $\|\eta_{t}\|_1=1$ for all $t$. For SEM, using \eqref{eq:dual} it follows that $\E\left[\|\eta_t\|_1^2\right]=O(1)$. In conclusion, we see that, for an absolute constant $C>0$, for all three models one has
		\begin{equation}
			\sum_{x}\E[E_t^{\rm nd}(x)]
				\le\frac{C}{n^2}+\left(1-\frac{2-4\E[X^2]}n\right)\sum_{x}\E[E_{t-1}^{\rm nd}(x)]\,.
		\end{equation}
		Iterating,	and using $\sum_xE_{0}^{\rm nd}	(x)=0$, 
		\begin{equation}\label{eq:zeta-B4}
			\sum_{x}\E[E_t^{\rm nd}]\le \frac{C}{n^2}\sum_{s=0}^\infty\left(1-\frac{2-4\E[X^2]}n\right)^s= \frac{C}{n(2-4\E[X^2])}=O\bigg(\frac1n\bigg)\,, 
		\end{equation}
		where $2-4\E[X^2]>0$ holds  by the non-degeneracy assumption. Remark that  the obtained estimate is uniform in $t\ge0$. By \eqref{eq:zeta-B2}, \eqref{eq:Ed-End}, \eqref{eq:zeta-B3} and \eqref{eq:zeta-B4},  we get \eqref{eq:topros}, as desired.
	\end{proof}

\end{document}
